\subsection{Monotone and bounded sequences}

The previous chapter introduced monotonicity and boundedness before limits. We can now explain why that combination is so important.

\begin{theorembox}[Monotone convergence theorem]
Every increasing sequence that is bounded above converges. Every decreasing sequence that is bounded below converges.
\end{theorembox}

For an increasing bounded sequence, the limit is the least upper bound of the set of its terms. For a decreasing bounded sequence, it is the greatest lower bound.

\begin{remarkbox}[Proof status]
The theorem rests on completeness of $\mathbb R$: every nonempty set bounded above has a least upper bound. The text explains how that boundary becomes the limit, but does not develop the completeness axiom itself.
\end{remarkbox}

\begin{examplebox}[A geometric sequence]
Let
\[
a_n=1-\frac1{2^n}.
\]
The sequence is increasing because $2^{-n}$ decreases. It is bounded above by $1$. Therefore it converges. Since $2^{-n}\to0$, its limit is
\[
\lim_{n\to\infty}\left(1-\frac1{2^n}\right)=1.
\]
\end{examplebox}

\begin{examplebox}[Why one condition is not enough]
The sequence $a_n=n$ is increasing but not bounded above, and it does not converge to a real number. The sequence $b_n=(-1)^n$ is bounded but not monotone, and it also does not converge.
\end{examplebox}

\begin{interpretationbox}[Direction plus confinement]
Monotonicity prevents repeated reversal of direction. Boundedness prevents escape. Completeness supplies the boundary value toward which the terms are forced.
\end{interpretationbox}
